% !TeX spellcheck = ru_RU
% !TEX root = vkr.tex

\section{Результаты экспериментов и их анализ} \label{sec:experiment}

В данном разделе приводятся результаты экспериментальной оценки 12~моделей на двух независимых бенчмарках: \texttt{benchmark-combined} (500~примеров: 326~GEO / 174~TAG) и \texttt{benchmark-real} (250~примеров: 203~GEO / 47~TAG). Методология подготовки данных и проведения экспериментов описана в разделе~\ref{sec:method}; определения метрик --- в разделе~\ref{sec:background}.

%%% ============================================================
\subsection{Оценка базовой модели} \label{subsec:baseline}
%%% ============================================================

Базовой линией (baseline) служит модель YandexGPT~5.1~Pro без дообучения (\texttt{pro\_baseline}). Результаты представлены в таблице~\ref{tab:baseline_results}.

\begin{table}[H]
\centering
\caption{Метрики базовой модели \texttt{pro\_baseline} на двух бенчмарках}
\label{tab:baseline_results}
\def\arraystretch{1.15}
\begin{tabular}{l r r r r}
    \toprule
    \textbf{Бенчмарк} & \textbf{Accuracy} & \textbf{Macro~F1} & \textbf{Ошибки} & \textbf{No Pred} \\
    \midrule
    combined (500) & $0{,}818$ & $0{,}526$ & 57 & 34 \\
    real (250)     & $0{,}828$ & $0{,}617$ & 28 & 15 \\
    \bottomrule
\end{tabular}
\end{table}

Низкое значение Macro~$F_1 = 0{,}526$ при относительно высокой accuracy $0{,}818$ указывает на выраженный дисбаланс качества между классами. Анализ матрицы ошибок выявляет систематический паттерн: все 57~ошибок на \texttt{benchmark\-combined} представляют собой ошибки типа TAG$\to$GEO, то есть запросы класса TAG массово классифицируются как GEO. Ни одного обратного случая (GEO$\to$TAG) не зафиксировано.

Данный bias объясняется тем, что базовая модель не дообучена для данного набора инструментов и по умолчанию предпочитает вызов \texttt{ResolveEntities} (GEO) как более общий. Помимо ошибок классификации, для 34~запросов (6{,}8\%) модель не генерирует \texttt{tool\_call} вовсе ($\rho_\varnothing = 6{,}8\%$), что является критическим сбоем в контексте AI-агента.

%%% ============================================================
\subsection{Основные результаты дообучения} \label{subsec:main_results}
%%% ============================================================

Результаты всех 12~моделей на \texttt{benchmark-combined} приведены в таблице~\ref{tab:combined_results}. Модели различаются по типу и объёму обучающих данных.

\begin{table}[H]
\centering
\caption{Результаты на \texttt{benchmark-combined} (500 примеров: 326~GEO / 174~TAG)}
\label{tab:combined_results}
\def\arraystretch{1.15}
\setlength\tabcolsep{0.3em}
{\small
\begin{tabular}{l l r r r r r}
    \toprule
    \textbf{Модель} & \textbf{Обуч. данные} & \textbf{Acc} & \textbf{M-F1} & \textbf{W-F1} & \textbf{Ош.} & \textbf{NP} \\
    \midrule
    pro\_baseline       & --- (без тюна)    & .818 & .526 & .829 & 57 & 34 \\
    \midrule
    lite\_latest\_v1    & Синт.\,${\sim}$1500 & .928 & .925 & .930 & 34 &  2 \\
    lite\_latest\_v2    & Синт.\,${\sim}$1500 & .984 & .985 & .986 &  6 &  2 \\
    lite\_latest\_v3    & Синт.\,${\sim}$1500 & .966 & .966 & .968 & 15 &  2 \\
    \midrule
    lite\_rc\_v1        & Синт.\,${\sim}$1500 & .980 & .981 & .982 &  8 &  2 \\
    lite\_rc\_v2        & Синт.\,${\sim}$1500 & .974 & .974 & .976 & 11 &  2 \\
    \midrule
    train\_synth\_50    & Синт.\,50        & .946 & .631 & .950 & 25 &  2 \\
    train\_synth\_250   & Синт.\,250       & .946 & .945 & .948 & 25 &  2 \\
    train\_synth\_500   & Синт.\,500       & .978 & .979 & .980 &  9 &  2 \\
    train\_synth\_1000  & Синт.\,908       & .956 & .955 & .958 & 20 &  2 \\
    \midrule
    train\_real\_250    & Реал.\,250         & .952 & .949 & .953 & 22 &  2 \\
    train\_mixed\_500   & Микс\,500          & .994 & .996 & .996 &  1 &  2 \\
    \bottomrule
\end{tabular}
}
\end{table}

Ключевые наблюдения:
\begin{enumerate}
    \item \textbf{Скачок при минимальном объёме данных.} Даже 50~синтетических примеров повышают accuracy с $0{,}818$ до $0{,}946$ ($+12{,}8$ п.п.), подтверждая гипотезу~\ref{hyp:lora}: даже минимальный объём примеров радикально улучшает качество, что согласуется с гипотезой поверхностного выравнивания~\cite{zhou2023lima}.

    \item \textbf{Устранение отказов.} Все дообученные модели снижают $\rho_\varnothing$ с $6{,}8\%$ (34~из~500) до $0{,}4\%$ (2~из~500).

    \item \textbf{Абсолютный лидер.} Модель \texttt{train\_mixed\_500}, обученная на сбалансированной комбинации реальных и синтетических данных (500~примеров), достигает accuracy~$0{,}994$ и Macro~$F_1 = 0{,}996$ с единственной ошибкой классификации.
\end{enumerate}

%%% ============================================================
\subsection{Валидация на реальных данных} \label{subsec:real_validation}
%%% ============================================================

Для оценки обобщающей способности моделей используется независимый бенчмарк \texttt{benchmark-real}, составленный из реальных пользовательских запросов. Результаты и разность с \texttt{benchmark-combined} приведены в таблице~\ref{tab:real_results}.

\begin{table}[H]
\centering
\caption{Результаты на \texttt{benchmark-real} (250 примеров: 203~GEO / 47~TAG)}
\label{tab:real_results}
\def\arraystretch{1.15}
\setlength\tabcolsep{0.35em}
\begin{tabular}{l r r r r r}
    \toprule
    \textbf{Модель} & \textbf{Acc} & \textbf{Macro~F1} & \textbf{Ош.} & \textbf{NP} & $\boldsymbol{\Delta}$\textbf{Acc} \\
    \midrule
    pro\_baseline       & $0{,}828$ & $0{,}617$ & 28 & 15 & $+0{,}010$ \\
    \midrule
    lite\_latest\_v1    & $0{,}904$ & $0{,}872$ & 22 &  2 & $-0{,}024$ \\
    lite\_latest\_v2    & $0{,}968$ & $0{,}959$ &  6 &  2 & $-0{,}016$ \\
    lite\_latest\_v3    & $0{,}932$ & $0{,}906$ & 15 &  2 & $-0{,}034$ \\
    \midrule
    lite\_rc\_v1        & $0{,}948$ & $0{,}921$ & 11 &  2 & $-0{,}032$ \\
    lite\_rc\_v2        & $0{,}944$ & $0{,}916$ & 12 &  2 & $-0{,}030$ \\
    \midrule
    train\_synth\_50    & $0{,}904$ & $0{,}584$ & 22 &  2 & $-0{,}042$ \\
    train\_synth\_250   & $0{,}892$ & $0{,}856$ & 25 &  2 & $-0{,}054$ \\
    train\_synth\_500   & $0{,}956$ & $0{,}942$ &  9 &  2 & $-0{,}022$ \\
    train\_synth\_1000  & $0{,}912$ & $0{,}882$ & 20 &  2 & $-0{,}044$ \\
    \midrule
    train\_real\_250    & $0{,}980$ & $0{,}978$ &  3 &  2 & $+0{,}028$ \\
    train\_mixed\_500   & $0{,}988$ & $0{,}991$ &  1 &  2 & $-0{,}006$ \\
    \bottomrule
\end{tabular}
\end{table}

Столбец $\Delta$Acc показывает разность accuracy на \texttt{benchmark-real} и \texttt{benchmark-combined}. Модели, обученные исключительно на синтетических данных, демонстрируют устойчивое снижение accuracy на реальных данных в диапазоне $-2{,}2\ldots-5{,}4$ п.п. Это свидетельствует об ограниченной обобщающей способности синтетических данных: модели частично переобучаются на характерные паттерны генерации, отсутствующие в реальных запросах.

Модель \texttt{train\_real\_250}, обученная на 250 реальных примерах, напротив, демонстрирует прирост $+2{,}8$ п.п.\ при переходе на реальный бенчмарк, что объяснимо совпадением распределений обучающей и тестовой выборок.

Модель \texttt{train\_mixed\_500} показывает минимальную деградацию: $\Delta\text{Acc} = -0{,}006$, что свидетельствует о высокой робастности, достигаемой за счёт сочетания разнородных обучающих данных.

%%% ============================================================
\subsection{Анализ кривой обучения} \label{subsec:learning_curve}
%%% ============================================================

Для исследования зависимости качества от объёма обучающих данных (RQ2) обучена серия моделей на синтетических выборках объёмом 50, 250, 500 и 908~примеров. Результаты на \texttt{benchmark-combined} представлены в таблице~\ref{tab:learning_curve}.

\begin{table}[H]
\centering
\caption{Кривая обучения: зависимость метрик от объёма синтетических данных}
\label{tab:learning_curve}
\def\arraystretch{1.15}
\begin{tabular}{r r r}
    \toprule
    \textbf{Объём} & \textbf{Accuracy} & \textbf{Macro~F1} \\
    \midrule
    50  & $0{,}946$ & $0{,}631$ \\
    250 & $0{,}946$ & $0{,}945$ \\
    500 & $0{,}978$ & $0{,}979$ \\
    908 & $0{,}956$ & $0{,}955$ \\
    \bottomrule
\end{tabular}
\end{table}

Кривая обучения обнаруживает нетривиальный немонотонный характер:
\begin{enumerate}
    \item \textbf{Насыщение accuracy при малых объёмах.} При переходе от 50 к 250 примерам accuracy не изменяется ($0{,}946$), однако Macro~$F_1$ скачкообразно возрастает с $0{,}631$ до $0{,}945$. Это означает, что 50 примеров достаточно для устранения основного bias базовой модели, но недостаточно для сбалансированного распознавания обоих классов.

    \item \textbf{Пик на 500 примерах.} Максимальное качество среди синтетических моделей достигается при 500~примерах: accuracy~$= 0{,}978$, Macro~$F_1 = 0{,}979$.

    \item \textbf{Деградация при увеличении до 908 примеров.} Дальнейшее наращивание объёма приводит к снижению accuracy до~$0{,}956$ ($-2{,}2$~п.п.). Такой эффект объяснятся переобучением на артефакты синтетической генерации: с ростом объёма выборки доля повторяющихся шаблонов увеличивается, и модель начинает эксплуатировать статистические закономерности, специфичные для синтетического распределения.
\end{enumerate}

%%% ============================================================
\subsection{Анализ типа обучающих данных} \label{subsec:data_type}
%%% ============================================================

Для проверки гипотезы~\ref{hyp:data_quality} сравним модели, обученные на данных различного типа при сопоставимых объёмах. Результаты сведены в таблице~\ref{tab:data_type_comparison}.

\begin{table}[H]
\centering
\caption{Сравнение моделей по типу обучающих данных}
\label{tab:data_type_comparison}
\def\arraystretch{1.15}
\setlength\tabcolsep{0.35em}
\begin{tabular}{l l r r r r}
    \toprule
    \textbf{Модель} & \textbf{Тип данных} & \multicolumn{2}{c}{\textbf{combined}} & \multicolumn{2}{c}{\textbf{real}} \\
    \cmidrule(lr){3-4} \cmidrule(lr){5-6}
     & & Acc & M-F1 & Acc & M-F1 \\
    \midrule
    train\_synth\_500  & Синтет.\,500  & $0{,}978$ & $0{,}979$ & $0{,}956$ & $0{,}942$ \\
    train\_real\_250   & Реал.\,250    & $0{,}952$ & $0{,}949$ & $0{,}980$ & $0{,}978$ \\
    train\_mixed\_500  & Микс\,500     & $0{,}994$ & $0{,}996$ & $0{,}988$ & $0{,}991$ \\
    \bottomrule
\end{tabular}
\end{table}

Наблюдается отчётливая специализация моделей по типу обучающих данных:
\begin{itemize}
    \item Модель на~синтетике (\texttt{train\_synth\_500}) превосходит модель на~реальных данных (\texttt{train\_real\_250}) на~смешанном бенчмарке ($+2{,}6$~п.п.), но~уступает ей на~реальном ($-2{,}4$~п.п.).
    \item \texttt{train\_real\_250} показывает лучший результат на \texttt{benchmark-real} ($0{,}980$), но её accuracy на \texttt{benchmark-combined} ниже ($0{,}952$), поскольку синтетические примеры в тестовой выборке имеют иную структуру.
    \item \texttt{train\_mixed\_500} объединяет преимущества обоих типов данных и доминирует на обоих бенчмарках: $0{,}994$ и $0{,}988$ соответственно.
\end{itemize}

Данный результат подтверждает гипотезу~\ref{hyp:data_quality}: качество (разнородность и репрезентативность) обучающих данных оказывает большее влияние на итоговые метрики, чем их объём. Модель \texttt{train\_mixed\_500}, обученная на 500~примерах, превосходит все модели, обученные на ${\sim}\,1500$ синтетических примерах.

%%% ============================================================
\subsection{Анализ ошибок} \label{subsec:error_analysis}
%%% ============================================================

\subsubsection{Паттерны ошибок классификации}

Для каждой модели определены направления ошибок: GEO$\to$TAG (запрос GEO ошибочно классифицирован как TAG) и TAG$\to$GEO (запрос TAG ошибочно классифицирован как GEO). Результаты представлены в таблице~\ref{tab:error_patterns}.

\begin{table}[H]
\centering
\caption{Паттерны ошибок классификации на \texttt{benchmark-combined}}
\label{tab:error_patterns}
\def\arraystretch{1.15}
\setlength\tabcolsep{0.4em}
\begin{tabular}{l r r l}
    \toprule
    \textbf{Модель} & \textbf{G$\to$T} & \textbf{T$\to$G} & \textbf{Домин.\,ошибка} \\
    \midrule
    pro\_baseline      &  0 & 57 & TAG массово как GEO \\
    train\_synth\_*    &  8--24 &  1 & GEO массово как TAG \\
    train\_real\_250   &  1 & 21 & TAG$\to$GEO (дисбал.) \\
    train\_mixed\_500  &  0 &  1 & Практически нет \\
    \bottomrule
\end{tabular}
\end{table}

Обнаруживается характерная инверсия bias: базовая модель систематически классифицирует TAG-запросы как GEO, тогда как дообученные на синтетике модели совершают преимущественно обратную ошибку (GEO$\to$TAG). Это объясняется тем, что синтетические обучающие данные перепредставляют класс TAG, смещая решающую границу модели.

Модель \texttt{train\_mixed\_500} практически свободна от обоих типов bias (0~ошибок GEO$\to$TAG, 1~ошибка TAG$\to$GEO), что подтверждает эффективность сбалансированного микса данных для устранения систематических смещений.

\subsubsection{Консервация ошибок}

Анализ пересечения ошибок на двух бенчмарках (\texttt{combined} и \texttt{real}) выявляет важный феномен: многие ошибки воспроизводимы, что свидетельствует об их детерминированной природе (таблица~\ref{tab:error_conservation}).

\begin{table}[H]
\centering
\caption{Консервация ошибок: совпадение ошибок на двух бенчмарках (combined / real)}
\label{tab:error_conservation}
\def\arraystretch{1.15}
\begin{tabular}{l c c}
    \toprule
    \textbf{Модель} & \textbf{GEO$\to$TAG} & \textbf{TAG$\to$GEO} \\
    \midrule
    lite\_latest\_v2   & 5\,/\,5 & 1\,/\,1 \\
    train\_synth\_500  & 8\,/\,8 & 1\,/\,1 \\
    train\_mixed\_500  & 0\,/\,0 & 1\,/\,1 \\
    \bottomrule
\end{tabular}
\end{table}

Для модели \texttt{lite\_latest\_v2} все 5~ошибок типа GEO$\to$TAG и 1~ошибка TAG$\to$GEO на \texttt{benchmark-combined} совпадают с ошибками на \texttt{benchmark-real}. Аналогичная картина наблюдается для \texttt{train\_synth\_500}. Это означает, что ошибки обусловлены не стохастикой генерации, а систематическими слабостями модели на определённых типах запросов (<<трудные>> примеры).

%%% ============================================================
\subsection{Статистическая значимость} \label{subsec:significance}
%%% ============================================================

\subsubsection{Критерий Кохрена}

Для проверки глобальной гипотезы о равенстве долей ошибок всех 12~моделей применён $Q$-критерий Кохрена (Cochran's~$Q$~test), являющийся обобщением критерия Макнемара на случай $k > 2$ связанных выборок (раздел~\ref{sec:background}):

\begin{itemize}
    \item \texttt{benchmark-combined}: $Q = 348{,}0$, $p < 10^{-6}$;
    \item \texttt{benchmark-real}: $Q = 128{,}8$, $p < 10^{-6}$.
\end{itemize}

На обоих бенчмарках нулевая гипотеза $H_0\colon \pi_1 = \pi_2 = \ldots = \pi_{12}$ отвергается при любом разумном уровне значимости. Следовательно, модели статистически значимо различаются по доле ошибок.

\subsubsection{Попарные сравнения (критерий Макнемара)}

Для попарного сравнения моделей применён точный критерий Макнемара с поправкой Бонферрони на множественные сравнения. При 12~моделях число пар составляет $\binom{12}{2} = 66$, скорректированный уровень значимости $\alpha_{\mathrm{adj}} = 0{,}05 / 66 \approx 0{,}0008$.

\paragraph{\texttt{benchmark-combined}.}
\begin{itemize}
    \item \texttt{train\_mixed\_500} статистически значимо превосходит большинство моделей ($p < \alpha_{\mathrm{adj}}$ для всех пар, кроме \texttt{lite\_latest\_v2}).
    \item Пара \texttt{lite\_latest\_v2} --- \texttt{lite\_rc\_v1} неразличима: $p = 0{,}73 \gg \alpha_{\mathrm{adj}}$.
\end{itemize}

\paragraph{\texttt{benchmark-real}.}
\begin{itemize}
    \item \texttt{train\_mixed\_500} и \texttt{train\_real\_250} статистически неразличимы: $p = 0{,}50$.
    \item Обе модели статистически значимо превосходят все модели, обученные исключительно на синтетических данных.
\end{itemize}

Данные результаты означают, что на реальных данных наличие реальных примеров в обучающей выборке является ключевым фактором: модель \texttt{train\_real\_250} (250~реальных примеров) статистически неотличима от \texttt{train\_mixed\_500} (500~смешанных), тогда как \texttt{train\_synth\_500} (500~синтетических) значимо уступает обеим.

\subsubsection{Согласованность предсказаний}

Анализ согласованности позволяет оценить долю примеров, на которых все 12~моделей дают одинаковый ответ:
\begin{itemize}
    \item \texttt{benchmark-combined}: $368$ из $500$ примеров ($73{,}6\%$) классифицированы правильно всеми 12~моделями; $0$~примеров ошибочны у всех~12.
    \item \texttt{benchmark-real}: $175$ из $250$ примеров ($70{,}0\%$) классифицированы правильно всеми 12~моделями.
\end{itemize}

Отсутствие примеров, ошибочных у всех 12~моделей, означает, что каждый запрос в бенчмарке в принципе решаем хотя бы одной из рассмотренных моделей. Оставшиеся ${\sim}\,27\text{--}30\%$ примеров составляют <<зону разногласий>>, где модели расходятся в предсказаниях, и именно на этих примерах сосредоточены различия в качестве.

%%% ============================================================
\subsection{Ответы на исследовательские вопросы} \label{subsec:rq_answers}
%%% ============================================================

\paragraph{\textsc{RQ1}: Каково влияние дообучения на метрики классификации?}

Дообучение методом LoRA приводит к статистически значимому улучшению всех метрик. Лучшая модель \texttt{train\_mixed\_500} повышает accuracy с $0{,}818$ до $0{,}994$ ($+17{,}6$~п.п.) и Macro~$F_1$ с $0{,}526$ до $0{,}996$ ($+0{,}470$). Даже наименее успешная дообученная модель (\texttt{lite\_latest\_v1}, accuracy~$= 0{,}928$) превосходит baseline на $+11{,}0$~п.п. Число отказов ($\rho_\varnothing$) снижается с $6{,}8\%$ до $0{,}4\%$ для всех дообученных моделей. Результат статистически значим по критерию Кохрена ($Q = 348{,}0$, $p < 10^{-6}$) и попарным критериям Макнемара.

\paragraph{\textsc{RQ2}: Как зависит качество от объёма обучающих данных?}

Зависимость немонотонна. При фиксированном типе данных (синтетика) accuracy растёт от $0{,}946$ (50~примеров) до пика $0{,}978$ (500~примеров), после чего снижается до~$0{,}956$ (908~примеров). Наиболее информативна динамика Macro~$F_1$: скачок с $0{,}631$ до $0{,}945$ при переходе от 50 к 250~примерам указывает на то, что именно этот диапазон объёмов критичен для достижения сбалансированного распознавания обоих классов. Деградация при 908~примерах свидетельствует о переобучении на синтетическое распределение.

\paragraph{\textsc{RQ3}: Как тип обучающих данных влияет на робастность модели?}

Тип данных оказывает определяющее влияние. Модели на чистой синтетике демонстрируют деградацию $-2{,}2\ldots -5{,}4$~п.п.\ при переходе на реальные данные. Модель \texttt{train\_real\_250} показывает обратный эффект: $+2{,}8$~п.п.\ на \texttt{benchmark-real} по сравнению с \texttt{benchmark-combined}. Оптимальная стратегия~--- смешанный датасет: \texttt{train\_mixed\_500} достигает accuracy~$\geq 0{,}988$ на обоих бенчмарках с деградацией всего $-0{,}6$~п.п. Это подтверждает гипотезу~\ref{hyp:data_quality}.

\paragraph{\textsc{RQ4}: Статистически значимы ли различия между моделями?}

Да. Критерий Кохрена отвергает гипотезу о равенстве моделей на обоих бенчмарках ($p < 10^{-6}$). Попарные сравнения Макнемара с поправкой Бонферрони ($\alpha_{\mathrm{adj}} \approx 0{,}0008$) показывают, что \texttt{train\_mixed\_500} значимо превосходит большинство моделей на \texttt{benchmark-combined}. На \texttt{benchmark-real} модели \texttt{train\_mixed\_500} и \texttt{train\_real\_250} статистически неразличимы ($p = 0{,}50$) и обе значимо лучше всех чисто синтетических моделей. Пара \texttt{lite\_latest\_v2}~--- \texttt{lite\_rc\_v1} также неразличима ($p = 0{,}73$).

\clearpage
